%% 
%% 版权所有 2019-2024 Elsevier Ltd
%% 
%% 本文件是 CAS Bundle 的一部分。
%% --------------------------------------
%% 
%% 本文件可在 LaTeX Project Public License 条款下发布和/或修改；
%% 许可证版本可为 1.3c，或你选择的任何更新版本。
%% 最新版本见：
%%    http://www.latex-project.org/lppl.txt
%% LaTeX 1999/12/01 或之后版本的标准发行版中已包含 1.3c 或更新版本。
%% 
%% CAS Bundle 所有文件的列表见 manifest.txt。
%% 
%% cas-dc 文档类的文章模板：
%% 用于双栏排版。

\documentclass[a4paper,fleqn]{cas-dc}

% 如果首页信息（标题、作者、摘要等）超过一页，
% 请使用 longmktitle 选项。

%\documentclass[a4paper,fleqn,longmktitle]{cas-dc}

%\usepackage[numbers]{natbib}
%\usepackage[authoryear]{natbib}
\usepackage[authoryear,longnamesfirst]{natbib}

%%% 作者自定义宏
\def\tsc#1{\csdef{#1}{\textsc{\lowercase{#1}}\xspace}}
\tsc{WGM}
\tsc{QE}
%%%

% 如有需要，可取消注释并使用
%\newtheorem{theorem}{Theorem}
%\newtheorem{lemma}[theorem]{Lemma}
%\newdefinition{rmk}{Remark}
%\newproof{pf}{Proof}
%\newproof{pot}{Proof of Theorem \ref{thm}}

\begin{document}
\let\WriteBookmarks\relax
\def\floatpagepagefraction{1}
\def\textpagefraction{.001}

% 短标题
\shorttitle{}    

% 短作者列表
\shortauthors{}  

% 论文主标题
\title [mode = title]{}  

% 标题脚注标记
% 示例：\tnotemark[1]
\tnotemark[1] 

% 标题脚注 1
% 示例：\tnotetext[1]{标题脚注内容}
\tnotetext[1]{} 

% 第一作者
%
% 可选参数：如有需要可使用
% 示例：\author[1,3]{Author Name}[type=editor,
%       style=chinese,
%       auid=000,
%       bioid=1,
%       prefix=Sir,
%       orcid=0000-0000-0000-0000,
%       facebook=<facebook id>,
%       twitter=<twitter id>,
%       linkedin=<linkedin id>,
%       gplus=<gplus id>]

\author[1]{}%[<options>]

% 通讯作者标记
\cormark[1]

% 第一作者脚注标记
\fnmark[1]

% 第一作者邮箱
\ead{}

% 第一作者个人主页或 URL
\ead[url]{}

% CRediT 作者贡献声明
% 示例：\credit{Conceptualization of this study, Methodology, Software}
\credit{}

% 地址/作者单位
\affiliation[1]{organization={},
            addressline={}, 
            city={},
%          citysep={}, % 如果城市和邮编之间不需要逗号，可取消注释
            postcode={}, 
            state={},
            country={}}

\author[2]{}%[]

% 第二作者脚注标记
\fnmark[2]

% 第二作者邮箱
\ead{}

% 第二作者个人主页或 URL
\ead[url]{}

% CRediT 作者贡献声明
\credit{}

% 地址/作者单位
\affiliation[2]{organization={},
            addressline={}, 
            city={},
%          citysep={}, % 如果城市和邮编之间不需要逗号，可取消注释
            postcode={}, 
            state={},
            country={}}

% 通讯作者说明文字
\cortext[1]{通讯作者}

% 作者脚注文本
\fntext[1]{}

% 无编号/无标记的标题说明
%\nonumnote{}

% 摘要
\begin{abstract}
在此填写摘要。\nocite{*}%% 从正式稿件中删除这一行。
\end{abstract}

% 如有图文摘要，可使用
%\begin{graphicalabstract}
%\includegraphics{}
%\end{graphicalabstract}

% 研究亮点
\begin{highlights}
\item 
\item 
\item 
\end{highlights}


%\nocite{*}

% 关键词
% 每个关键词之间用 \sep 分隔
\begin{keywords}
 \sep \sep \sep
\end{keywords}

\maketitle

% 正文
\section{}\label{}

% 有序列表
% 可在方括号中指定编号样式。
% 如果不指定，则使用默认样式。
%\begin{enumerate}[a)]
%\item 
%\item 
%\item 
%\end{enumerate}  

% 无序列表
%\begin{itemize}
%\item 
%\item 
%\item 
%\end{itemize}  

% 描述列表
%\begin{description}
%\item[]
%\item[] 
%\item[] 
%\end{description}  

\clearpage %% 从正式稿件中删除这一行

% 图
\begin{figure}%[]
  \centering
%    \includegraphics{}
    \caption{}\label{fig1}
\end{figure}


\begin{table}%[]
\caption{}\label{tbl1}
\begin{tabular*}{\tblwidth}{@{}LL@{}}
\toprule
  &  \\ % 表头行
\midrule
 & \\
 & \\
 & \\
 & \\
\bottomrule
\end{tabular*}
\end{table}

% 如有需要，可取消注释并使用
%\begin{theorem} 
%\end{theorem}

% 如有需要，可取消注释并使用
%\begin{lemma} 
%\end{lemma}

%% 附录部分以 \appendix 命令开始；
%% 之后的附录章节仍按普通 section 编写。
%% \appendix

\section{}\label{}

% 打印 CRediT 作者贡献明细
\printcredits

%% 加载参考文献样式文件
%\bibliographystyle{model1-num-names}
\bibliographystyle{cas-model2-names}

% 加载参考文献数据库
\bibliography{cas-refs}

% 作者简介
%\bio{}
% 在此填写作者简介。
%\endbio

%\bio{pic1}
% 在此填写带照片的作者简介。
%\endbio

\end{document}

